% Part: computability
% Chapter: computability-theory
% Section: complement-ce

\documentclass[../../../include/open-logic-section]{subfiles}

\begin{document}

\olfileid{cmp}{thy}{cmp}
\olsection{पूरक संचाखाली संगणनक्षमपणे प्रगणनीय संच बंद नसतात}

$A$ संगणनक्षमपणे प्रगणनीय आहे असे समजा. मग~$A$ चा पूरक संच,
$\Complement{A} = \Nat \setminus A$, हाही नेहमी संगणनक्षमपणे
प्रगणनीय असतो का? पुढील प्रमेय आणि अनुप्रमेय दाखवतात की उत्तर
``नाही'' असे आहे.

\begin{thm}
\ollabel{thm:ce-comp}
$A$ हा नैसर्गिक संख्यांचा कोणताही संच असू द्या. $A$ आणि
$\Complement{A}$ हे दोन्ही संगणनक्षमपणे प्रगणनीय असतील
तेव्हाच $A$ संगणनक्षम असतो.
\end{thm}

\begin{proof}
एका दिशेची सिद्धता सोपी आहे: $A$ संगणनक्षम असेल, तर
$\Complement{A}$ देखील संगणनक्षम आहे ($\Char{A} = 1 \tsub
\Char{\Complement{A}}$). म्हणून दोन्ही संच संगणनक्षमपणे प्रगणनीय आहेत.

उलट दिशेने, $A$ आणि~$\Complement{A}$ हे दोन्ही संगणनक्षमपणे
प्रगणनीय आहेत असे समजा. $A$ हे~$\cfind{d}$ चे परिभाषाक्षेत्र
आणि $\Complement{A}$ हे~$\cfind{e}$ चे परिभाषाक्षेत्र असू द्या.
$h$ ची व्याख्या अशी करा:
\[
h(x) = \umin{s}{(T(d,x,s) \lor T(e,x,s))}.
\]
म्हणजे, आदान~$x$ वर $h$ हे~$\cfind{d}$ किंवा~$\cfind{e}$
यांचे थांबणारे संगणन शोधते. $x \in A$ असेल, तर पहिल्या
पर्यायात शोध यशस्वी होतो; $x \in
\Complement{A}$ असेल, तर दुसऱ्या पर्यायात. म्हणून $h$
हे सर्वत्र परिभाषित संगणनक्षम फलन आहे. आता प्रत्येक~$x$
साठी $x \in A$ असेल तेव्हाच $T(d, x, h(x))$, म्हणजे
या आदानावर $\cfind{d}$ परिभाषित असते. $T(d, x, h(x))$
हा संगणनक्षम संबंध असल्यामुळे $A$ संगणनक्षम आहे.
%READERNOTE{OLCMP-020 — गोठवलेल्या इंग्रजी सिद्धतेत A चे परिभाषाक्षेत्र cfind_d असे ठरवूनही अंतिम कसोटीत T(e,x,h(x)) आणि cfind_e लिहिले आहे; ती कसोटी प्रत्यक्षात A च्या पूरकाची आहे. येथे d वापरले आहे. पुढील अनौपचारिक स्पष्टीकरणातील e,f ही जोडीही d,e अशी दुरुस्त केली आहे.}
\end{proof}

\begin{explain}
अनौपचारिक संगणकीय भाषेत ही गोष्ट अधिक सोपी आहे: $A$
निर्णीत करण्यासाठी आदान $x$ वर $\cfind{d}$ आणि $\cfind{e}$
यांच्या थांबणाऱ्या संगणनांचा शोध घ्या. त्यांपैकी एक नक्की थांबते;
$\cfind{d}$ थांबले, तर $x$ हा~$A$ मध्ये आहे; अन्यथा $x$
हा~$\Complement{A}$ मध्ये आहे.
\end{explain}

\begin{cor}
\ollabel{cor:comp-k}
$\Complement{K_0}$ संगणनक्षमपणे प्रगणनीय नाही.
\end{cor}

\begin{proof}
$K_0$ संगणनक्षमपणे प्रगणनीय आहे, पण संगणनक्षम नाही, हे
आपल्याला माहीत आहे. $\Complement{K_0}$ संगणनक्षमपणे प्रगणनीय
असता, तर \olref{thm:ce-comp} मुळे $K_0$ संगणनक्षम ठरला असता;
हे \olref[ncp]{thm:K-0} ला विरोधी आहे.
\end{proof}

\end{document}
